% a mashup of hipstercv, friggeri and twenty cv
% https://www.latextemplates.com/template/twenty-seconds-resumecv
% https://www.latextemplates.com/template/friggeri-resume-cv

\documentclass[lighthipster]{simplehipstercv}
% available options are: darkhipster, lighthipster, pastel, allblack, grey, verylight, withoutsidebar
% withoutsidebar
\usepackage[utf8]{inputenc}
\usepackage[default]{raleway}
\usepackage[margin=1cm, a4paper]{geometry}


%------------------------------------------------------------------ Variablen

\newlength{\rightcolwidth}
\newlength{\leftcolwidth}
\setlength{\leftcolwidth}{0.23\textwidth}
\setlength{\rightcolwidth}{0.75\textwidth}

%------------------------------------------------------------------
\title{CV - Maia Ortiz}
\author{Maia Ortiz}
\date{Marzo 2026}

\pagestyle{empty}
\begin{document}


\thispagestyle{empty}
%-------------------------------------------------------------

\section*{Start}

\simpleheader{headercolour}{Maia}{Ortiz}{Estudiante de Ingeniería Industrial}{white}



%------------------------------------------------

% this has to be here so the paracols starts..
\subsection*{}
\vspace{4em}

\setlength{\columnsep}{1.5cm}
\columnratio{0.23}[0.75]
\begin{paracol}{2}
\hbadness5000
%\backgroundcolor{c[1]}[rgb]{1,1,0.8} % cream yellow for column-1 %\backgroundcolor{g}[rgb]{0.8,1,1} % \backgroundcolor{l}[rgb]{0,0,0.7} % dark blue for left margin

\paracolbackgroundoptions

% 0.9,0.9,0.9 -- 0.8,0.8,0.8


\footnotesize
{\setasidefontcolour
\flushright
\begin{center}
    \roundpic{Foto Mai.jpeg}
\end{center}

\bg{cvgreen}{white}{Sobre mí}\\[0.5em]

{\footnotesize
Estudiante de Ingeniería Industrial en la UNCuyo, con gran capacidad analítica y orientación al detalle. Me destaco por mi responsabilidad, puntualidad y habilidad para la planificación organizada de proyectos técnicos.}
\bigskip

\bg{cvgreen}{white}{Personal} \\[0.5em]
Maia Ortiz

20 años

Mendoza, Argentina

2612469048
\bigskip

\bg{cvgreen}{white}{Areas de especialización} \\[0.5em]

 ~•~ Termodinámica \\
 ~•~ Mecánica de Fluidos \\
 ~•~ Electrotecnia \\
 ~•~ Química Orgánica \\

\bigskip



\bigskip

\bg{cvgreen}{white}{Intereses}\\[0.5em]

~•~ Sistemas de Productividad Personal \\
~•~ Optimización de Procesos \\ 
~•~ Análisis de Datos Técnicos \\
~•~ Bienestar y Disciplina \\
\bigskip

%\bg{cvgreen}{white}{Interests}\\[0.5em]

%\texttt{R} ~/~ \texttt{Android} ~/~ \texttt{Linux}

%\texttt{R} ~/~ \texttt{Android} ~/~ \texttt{Linux}

%\texttt{R} ~/~ \texttt{Android} ~/~ \texttt{Linux}

\vspace{4em}

\infobubble{\faAt}{cvgreen}{white}{maialortiz@gmail.com}
%\infobubble{\faTwitter}{cvgreen}{white}{-}
%\infobubble{\faFacebook}{cvgreen}{white}{-}
%\infobubble{\faGithub}{cvgreen}{white}{-}

\phantom{turn the page}

\phantom{turn the page}
}
%-----------------------------------------------------------
\switchcolumn

\small
\section*{Formación académica}

\begin{tabular}{r| p{0.42\textwidth} c}
    \cvevent{2024-Presente}{Ingeniería Industriall}{Estudiante}{UNCuyo \color{cvred}}{Formación integral con énfasis en ciencias térmicas y de fluidos. Desarrollo de algoritmos en \textbf{MATLAB} y \textbf{Octave} para la resolución de sistemas complejos mediante métodos numéricos (Jacobi, Runge-Kutta). Aplicación de normativas de seguridad y precisión en laboratorios de Química Orgánica y Electrotecnia.}{UNCuyo.jpg} \hspace{1cm} \\
    \cvevent{2019--2023}{Bachiller en Ciencias Sociales y Humanidades}{Colegio DAD}{Mendoza \color{cvred}}{Desarrollo de sólidas capacidades de análisis crítico, redacción técnica y comunicación efectiva. Reconocida por mi enfoque en la \textbf{planificación organizada} y el cumplimiento detallado de objetivos académicos en proyectos interdisciplinarios.}{dad.jpg} \hspace{1cm} \\
\end{tabular}
\vspace{3em}

\begin{minipage}[t]{0.35\textwidth}
\section*{Proyectos y experiencia}
\begin{tabular}{r p{0.6\textwidth} c}
    \cvdegree{2025}{Métodos Numéricos}{Algoritmos en Octave}{UNCuyo \color{headerblue}}{}{UNCuyo.jpg} \\
    \cvdegree{2025}{Física II}{Óptica y Polarización}{UNCuyo \color{headerblue}}{}{UNCuyo.jpg} \\
    \cvdegree{2024}{Química Orgánica}{Ensayos de Laboratorio}{UNCuyo \color{headerblue}}{}{UNCuyo.jpg}
\end{tabular}
\end{minipage}\hfill
\begin{minipage}[t]{0.3\textwidth}
\section*{Habilidades Técnicas}
\begin{tabular}{r @{\hspace{0.5em}}l}
     \bg{skilllabelcolour}{iconcolour}{html} &  \barrule{0.4}{0.5em}{cvpurple}\\
     \bg{skilllabelcolour}{iconcolour}{\LaTeX} & \barrule{0.55}{0.5em}{cvgreen} \\
     \bg{skilllabelcolour}{iconcolour}{Octave} & \barrule{0.5}{0.5em}{cvpurple} \\
     \bg{skilllabelcolour}{iconcolour}{python} & \barrule{0.25}{0.5em}{cvpurple} \\
     \bg{skilllabelcolour}{iconcolour}{R} & \barrule{0.1}{0.5em}{cvpurple} \\
\end{tabular}
\end{minipage}
{2025}{Física II}{Óptica y Polarización}{UNCuyo}
\section*{Curriculum}
\begin{tabular}{r| p{0.5\textwidth} c}
    \cvevent{2025}{Métodos Numéricos}{Simulación Numérica}{UNCuyo \color{cvred}}{Implementación de algoritmos en Octave (Runge-Kutta y Jacobi) para la resolución de sistemas complejos}{UNCuyo.jpg} \\
    \cvevent{2025}{Física II}{Óptica y Polarización}{UNCuyo \color{cvred}}{Uso de instrumental de precisión para ensayos de laboratorio y análisis de datos experimentales}{UNCuyo.jpg} \\
\end{tabular}
\vspace{3em}

\begin{minipage}[t]{0.3\textwidth}
\section*{Logros Académicos}
\begin{tabular}{>{\footnotesize\bfseries}r >{\footnotesize}p{0.55\textwidth}}
    2026 & Presentación académica (UNCuyo) \\
    2025 & Certificación interna de Seguridad en Laboratorio \\
    2023 & Egresada con Honores - Bachiller en Cs. Sociales (DAD)
\end{tabular}
\bigskip

\section*{Idiomas}
\begin{tabular}{l | ll}
\textbf{Español} & Nativo & \pictofraction{\faCircle}{cvgreen}{4}{black!30}{1}{\tiny} \\
\textbf{Inglés} & Técnico & \pictofraction{\faCircle}{cvgreen}{3}{black!30}{1}{\tiny} \\
\textbf{Francés} & Básico & \pictofraction{\faCircle}{cvgreen}{1}{black!30}{3}{\tiny} 
\end{tabular}
\bigskip

\end{minipage}\hfill
\begin{minipage}[t]{0.3\textwidth}
\section*{Proyectos}
\begin{tabular}{>{\footnotesize\bfseries}r >{\footnotesize}p{0.7\textwidth}}
    2025 & \emph{Simulación en Octave}: Implementación de los métodos Jacobi y Runge-Kutta para la resolución de sistemas de ecuaciones y modelos dinámicos. \\
    2026 & Bitácora de Aprendizaje. \emph{Desarrollo de un perfil profesional detallado y registro de procesos técnicos en ingeniería.}
\end{tabular}
\bigskip

\section*{Presentaciones}
\begin{tabular}{>{\footnotesize\bfseries}r >{\footnotesize}p{0.6\textwidth}}
    Mar. 2026 & ``Estrategias de Trabajo en Equipo y Colaboración Académica'', en: \emph{Ingeniería Industrial} UNCuyo.
\end{tabular}
\end{minipage}






\vfill{} % Whitespace before final footer

%----------------------------------------------------------------------------------------
%	FINAL FOOTER
%----------------------------------------------------------------------------------------
\setlength{\parindent}{0pt}
\begin{minipage}[t]{\rightcolwidth}
\begin{center}\fontfamily{\sfdefault}\selectfont \color{black!70}
{\small Maia Ortiz \icon{\faEnvelopeO}{cvgreen}{} Mendoza \icon{\faMapMarker}{cvgreen}{} Argentina \icon{\faPhone}{cvgreen}{} 2612469048 \newline\icon{\faAt}{cvgreen}{} \protect\url{maialortiz@gmail.com}
}
\end{center}
\end{minipage}

\end{paracol}

\end{document}
